\section{Related work}
\label{sec:related}

\paragraph{Leakage.}
Leakage has been formalised for data mining~\citep{kaufman2012leakage} and catalogued as a cause of irreproducible results across applied fields~\citep{kapoor2023leakage,lones2024pitfalls}. Image-level instead of subject-level splits inflate accuracy for brain MRI~\citep{yagis2021mri} and retinal OCT~\citep{tampu2022oct}, and COVID-19 radiograph models suffered from many such pitfalls~\citep{roberts2021pitfalls}. \citet{gerz2025augmentation} show with two detection case studies that augmenting before splitting inflates mAP, and \citet{babu2024automotive} detect leakage between consecutive frames of driving video. These works measure leakage when the grouping is known. We study the case where it is not: we quantify why random splits of pre-augmented releases are nearly uninformative and can err in either direction, show that pair-level de-duplication does not remove the leak, and give a calibrated, metadata-free audit.

\paragraph{Duplicates in benchmarks and corpora.}
\citet{barz2020cifair} found near-duplicates of test images in the CIFAR training sets; \citet{recht2019imagenet} showed accuracy drops on freshly collected test sets; \citet{ramos2025leakage} measure exact and near-identical overlap within and across vision datasets; \citet{northcutt2021pervasive} audit test-set label errors. Web-scale de-duplication uses embeddings~\citep{abbas2023semdedup,webster2023laion,oquab2024dinov2} or hashing~\citep{lee2022dedup}, and copy detection learns edit-robust descriptors~\citep{pizzi2022sscd}. These target near-identical pairs at a fixed threshold. Pre-augmented releases differ in two ways that matter: copies are deliberately dissimilar, and removing leakage needs groups, not pairs. Our null turns the threshold into an error rate, using perceptual hashing~\citep{zauner2010phash}, ratio-test and local-scaling ideas~\citep{lowe2004sift,conneau2018csls} against hubs~\citep{radovanovic2010hubs}, and extreme-value theory~\citep{coles2001evt}.
